% Phụ lục F — bảng đối sánh giải pháp công nghệ, tách khỏi thân báo cáo.
% Phần kết luận và lý do chính nằm tại mục 5.2.

\section*{Phụ lục F. Đối sánh các hướng tiếp cận cho mô-đun AI}
\label{apdx:congnghe}
\addcontentsline{toc}{section}{Phụ lục F. Đối sánh các hướng tiếp cận cho mô-đun AI}

Bảng~\ref{tab:sosanh-ai} đối sánh các hướng tiếp cận cho từng mô-đun AI; đối sánh các
phương án công nghệ nền và kết luận lựa chọn nằm ở Bảng~\ref{tab:congnghe}, mục~\ref{sec:giaiphapcongnghe}.
Cột nhận xét là nhận định của nhóm dựa trên tài liệu công bố của từng phương án,
chưa phải số liệu đo; các phương án được chọn sẽ được đo thật ở Chương~\ref{chap:danhgia}.

\begingroup
\tblsetup\tblzebra
\begin{longtable}{L{2.60cm} L{3.74cm} L{4.67cm} L{3.53cm}}
\caption{Đối sánh các hướng tiếp cận cho từng mô-đun AI}
\label{tab:sosanh-ai} \\
\toprule
\rowcolor{tblheadbg}
\textbf{Mô-đun} & \textbf{Hướng tiếp cận} & \textbf{Nhận xét chính} & \textbf{Kết luận} \\
\midrule
\endfirsthead
\toprule
\rowcolor{tblheadbg}
\textbf{Mô-đun} & \textbf{Hướng tiếp cận} & \textbf{Nhận xét chính} & \textbf{Kết luận} \\
\midrule
\endhead
\bottomrule
\endfoot
Mô hình nhúng văn bản & BGE-M3 & Đa ngôn ngữ có tiếng Việt, dùng được cho truy vấn ngắn lẫn mô tả dài, mã nguồn mở & \textbf{Chọn} \\
& Multilingual-E5 & Đa ngôn ngữ, nhẹ hơn; dùng làm đối chứng khi đo nDCG@10 & Chọn làm đối chứng \\
& SBERT huấn luyện riêng cho tiếng Việt & Tốt với tiếng Việt nhưng thường cần tách từ trước và kém với truy vấn trộn tiếng Anh & Không chọn \\
Tìm kiếm theo ngữ nghĩa & Tìm kiếm lai: từ khóa và vector, hợp nhất theo thứ hạng & Giữ điểm mạnh của từ khóa với tên riêng; mô hình nhúng mã nguồn mở, không tốn phí gọi API & \textbf{Chọn} \\
& Tìm theo vector thuần & Dễ bỏ sót kết quả khớp chính xác tên sản phẩm, vùng trồng & Không chọn \\
& Mô hình ngôn ngữ lớn chuyển câu hỏi thành bộ lọc & Cần API trả phí hoặc hạ tầng chạy mô hình lớn; khó đánh giá bằng số & Không chọn \\
& Tìm kiếm dạng hội thoại & Thực chất là trợ lý ảo, đã bỏ khỏi phạm vi & Không chọn \\
Gợi ý sản phẩm & Collaborative filtering thuần & Không hoạt động khi chưa có lịch sử tương tác & Không chọn \\
& Content-based filtering thuần & Chạy được ngay nhưng gợi ý đơn điệu khi dữ liệu tăng & Không chọn \\
& Mô hình hybrid & Dùng đặc trưng nông sản khi chưa có dữ liệu, chuyển dần sang hành vi; phần hiện thực chỉ có sản phẩm tương tự theo vector & \textbf{Chọn} \\
& Mạng nơ-ron & Đòi hỏi lượng tương tác lớn hơn nhiều so với khả năng thu thập & Không chọn \\
Dự báo giá & Mô hình time series có mùa vụ & Phù hợp chuỗi time series ngắn nhưng có tính mùa vụ rõ; dễ giải thích & \textbf{Chọn} \\
& Gradient boosting trên đặc trưng thời gian & Khai thác được biến ngoại sinh; dùng làm mô hình đối chứng & Chọn làm đối chứng \\
& Mạng nơ-ron hồi quy (RNN, LSTM) & Cần lượng dữ liệu vượt xa khả năng thu thập trong đồ án & Không chọn \\
Giám sát bất thường & Luật ngưỡng cấu hình được & Kiểm chứng được, minh bạch, không cần dữ liệu huấn luyện & \textbf{Chọn} \\
& Unsupervised learning & Không đủ lưu lượng thật để học phân bố bình thường & Hướng phát triển \\
& Supervised learning & Không có tập nhãn gian lận đã xác nhận để đo chất lượng & Hướng phát triển \\
\end{longtable}
\endgroup
